\chapter{Prefácio}
Este \textit{template} foi reescrito pelo Prof. Dr. Jonatha Rodrigues da Costa (JRC) a partir do modelo original do IFCE \textit{campus} Fortaleza. Há orientações específicas para cada elemento do texto, fundamentadas em sua experiência como docente, pesquisador e orientador de Trabalhos de Conclusão de Curso na área de Engenharia.

Dessa forma, ao longo deste documento serão apresentadas recomendações metodológicas, técnicas e científicas que refletem as expectativas do orientador quanto à organização do trabalho, ao rigor metodológico, à profundidade técnica e à qualidade da escrita científica.
%

\section{Estruturação}
Orienta-se que o TCC seja escrito de forma clara, técnica e estruturada como segue:

\begin{enumerate}
	\item Introdução
	\begin{enumerate}
		\item Contextualização
		\item Objetivos
		\item Contribuições
		\item Publicações
		\item Organização do trabalho
	\end{enumerate}
	\item Estado da arte
	\begin{enumerate}
		\item Considerações iniciais
		\item Fronteira do conhecimento nos assuntos envolvidos
		\item Visão holística do estado da arte 
		\item Considerações finais
	\end{enumerate}
	\item Fundamentação teórica
	\begin{enumerate}
		\item Considerações iniciais
		\item Conteúdo seminal correlato ao assunto do TCC 
		\item Considerações finais
	\end{enumerate}	
	\item Solução proposta 
	\begin{enumerate}
		\item Arquitetura holística do sistema
		\item Modelo matemático do sistema
		\item Fluxograma do sistema
		\item Considerações finais
	\end{enumerate}
	\item Simulações, Análises e Discussão de resultados
	\begin{enumerate}
		\item Considerações iniciais
		\item Cenários, métodos e parâmetros de simulação/teste/coleta de dados
		\item Apresentação de discussão de resultados: tabelas, gráficos, quadros
		\item Considerações finais
	\end{enumerate}
	\item Conclusões e trabalhos futuros
	\item Referências
	\item Apêndices
	\item Anexos
\end{enumerate}


Esta estruturação garante um mínimo de qualidade técnico-científica a ser aplicada durante a produção de TCC do concluinte de engenharia.
Isto visa, sobre tudo, maximizar o tempo do discente e orientador, reduzindo riscos de ir à defesa com material insuficiente para tanto. Note que este documento está estruturado em \LaTeX\ contendo arquivos \texttt{.tex} já organizados conforme estrutura supracitada. Além disso, há pastas específicas destinadas a organizar o conteúdo em elementos \textbf{pré-textuais}, \textbf{textuais} e \textbf{pós-textuais}.
Desse forma, o concluinte redator do texto deve alterar os arquivos, ou apenas comentar aqueles não aplicáveis ao seu TCC.


% =========================================================
\section{Estrutura lógica e princípio orientador do TCC}
% =========================================================

Um TCC científico deve ser construído como uma sequência lógica e integrada,
na qual cada capítulo responde a uma questão fundamental da pesquisa:

\begin{center}
	\begin{tabular}{>{\bfseries}p{4cm}p{9cm}}
		\toprule
		\textbf{Etapa} & \textbf{Questão central} \\
		\midrule
		Problema & Por que o problema importa? \\
		Estado da Arte & O que já foi feito e o que ainda falta? \\
		Lacuna & O que permanece sem solução? \\
		Objetivo & O que este trabalho pretende realizar? \\
		Fundamentação & Que conhecimento sustenta a pesquisa? \\
		Solução & O que foi desenvolvido? \\
		Metodologia & Como a solução foi avaliada? \\
		Resultados & O que foi observado? \\
		Discussão & O que os resultados significam? \\
		Conclusão & O que foi efetivamente demonstrado? \\
		Trabalhos futuros & O que ainda precisa ser investigado? \\
		\bottomrule
	\end{tabular}
\end{center}

\textbf{
Note que este fluxo deve guiar também a produção da  apresentação (\textit{slildes}) para o defesa.
}

\textbf{Princípio orientador:} no corpo do texto, cada seção deve responder perguntas, cada afirmação relevante deve possuir
evidência, cada resultado deve ser interpretado e cada conclusão deve estar
sustentada pelos resultados apresentados. Essa estrutura garante que o TCC seja mais do que uma descrição de atividades:
ele deve constituir uma \textbf{argumentação científica coerente, fundamentada
	em evidências e capaz de demonstrar uma contribuição}.

Note que uma banca de avaliação pode sempre ser mais rigorosa do que aquilo que se imagina. Portanto, ações proativas mitigam frustrações.

A seguir serão apresentas perguntas-chave para cada capítulo do TCC.

